\documentclass{article}
\usepackage[UTF8]{ctex}
\usepackage{amsmath,amssymb,amsfonts}
\usepackage{enumitem}
\usepackage{geometry}
\geometry{a4paper, margin=1in}
% === 错题本排版优化（不参与题目标识） ===
\setlength{\parindent}{0pt}
\setlength{\parskip}{0.35em plus 0.10em minus 0.08em}
\setlength{\emergencystretch}{3em}
\setlist{itemsep=0.25em, topsep=0.40em, parsep=0pt}
\allowdisplaybreaks[2]
% === 错题本排版优化结束 ===
\begin{document}
\section*{西安电子科技大学2025年数学分析真题}
\section*{191．西安电子科技大学2025年数学分析试卷}

\subsection{题目1}
1． \(\lim _{x \rightarrow 0} \frac{1-(\cos x)^{\sin x}}{x^{3}}\) ．

\subsection{题目2}
2． \(\int \frac{1+\arcsin x}{\left(1-x^{2}\right)^{\frac{3}{2}}} \mathrm{~d} x\) ．

\subsection{题目3}
3．设 \(a_{n}=\int_{-1}^{1} \frac{x^{2 n}}{1+e^{x}} \mathrm{~d} x, n=0,1,2, \cdots\) ，求级数 \(\sum_{n=0}^{\infty}(-1)^{n} a_{n}\) 的和．

\subsection{题目4}
4．计算 \(I=\iint_{\sum} \sqrt{x^{2}+y^{2}} \mathrm{~d} S\) ，其中 \(\sum: x^{2}+y^{2}+z^{2}=1, z \geq 0\) ．

\subsection{题目5}
5．计算 \(\iint_{S} x^{3} \mathrm{~d} y \mathrm{~d} z, S\) 是椭球面 \(\frac{x^{2}}{a^{2}}+\frac{y^{2}}{b^{2}}+\frac{z^{2}}{c^{2}}=1\) 的上半部分，取上侧．

\subsection{题目6}
6．设 \(f(x)\) 在 \(a\) 的某邻域内具有连续的 \(n+1\) 阶导函数，求 \(\lim _{x \rightarrow a} \frac{\mathrm{~d}^{n}}{\mathrm{~d} x^{n}}\left[\frac{f(z)-f(a)}{z-a}\right]\) ．

二、判断下列命题是否正确，若正确请证明，错误则给出反例。每题7分，共28分。

\begin{itemize}

\subsection{题目7}
  \item[1．] 函数列 \(\left\{f_{n}(x)\right\}=\left\{\frac{1}{1+n^{2} x}\right\}\) 在（0．1）上一致收敛。

\subsection{题目8}
  \item[2．] 对任意正整数 \(n\) 和 \(p\) ，均有 \(\left|x_{n+p}-x_{n}\right|<\frac{p}{n^{2}}\) ，则数列（ \(x_{n}\) ）是柯西列。

\subsection{题目9}
  \item[3．] 设二元函数 \(f(x, y)\) 有二阶连续偏导数，\(\left(x_{0}, y_{0}\right)\) 是 \(f(x, y)\) 的严格极小值点，则
\[
f_{x x}\left(x_{0}, y_{0}\right) f_{y y}\left(x_{0}, y_{0}\right)-f_{x y}^{2}\left(x_{0}, y_{0}\right)>0 .
\]

\subsection{题目10}
  \item[4．] 设级数 \(\sum_{n=1}^{\infty} u_{n}\) 收敛，则级数 \(\sum_{n=1}^{\infty} u_{n}^{3}\) 收敛。
\end{itemize}

三．（13 分）有 \(p \in \mathbb{R}\) ，讨论反常积分 \(\int_{0}^{+\infty} \frac{\sin x}{x^{p}} \mathrm{~d} x\) 的敛散性，若收敛，判断是条件收敛还是绝对收敛．

四．（13 分）设函数

\[
k(x, y)= \begin{cases}y(1-x), & y \leq x \\ x(1-y), & y>x\end{cases}
\]

\(f(y)\) 在 \([0,1]\) 上连续．证明：\(u(x)=\int_{0}^{1} k(x, y) f(y) \mathrm{d} y\) 在 \([0,1]\) 上满足微分方程 \(\frac{\mathrm{d}^{2} u}{\mathrm{~d} x^{2}}+f(x)=0\) ．

五．（14 分）\(\varphi(y)\) 有连续的导函数，且对右半平面 \(x>0\) 内任意闭合曲线 \(C\) ，均有

\[
\oint_{C} \frac{\varphi(y) \mathrm{d} x+3 x y^{2} \mathrm{~d} y}{2 x^{2}+y^{6}}=0 .
\]

\begin{itemize}
  \item[（1）] 求 \(\varphi(y)\) 的表达式。
  \item[（2）] 设曲线 \(L: x^{2}+2 y^{2}=1\) ，取逆时针方向，求 \(\oint_{L} \frac{\varphi(y) \mathrm{d} x+3 x y^{2} \mathrm{~d} y}{2 x^{2}+y^{6}}\) ．
\end{itemize}

六．（13 分）设 \(f_{i}(x)(i=1,2, \cdots, n)\) 为 \([a, b]\) 上的连续函数，证明：

\[
\sum_{k=1}^{n}\left(\int_{a}^{b} f_{k}(x) \mathrm{d} x\right)^{2} \leq\left(\int_{a}^{b} \sqrt{\sum_{k=1}^{n} f_{k}^{2}(x)} \mathrm{d} x\right)^{2}
\]

七．（13 分）设向量值函数 \(f:(-1,1) \rightarrow \mathbb{R}^{2}\) 有连续导数，且 \(f(0)=0, f^{\prime}(0) \neq 0\) ，证明：存在 \(\varepsilon \in(0,1)\) ，使得 \(\|\mathbf{f}(t)\|\) 关于 \(t\) 在 \((0, \varepsilon)\) 上单调递增．

八．（14 分）设函数 \(f\) 在区间 \([a, b]\) 上有定义且处处存在极限，证明：

\begin{itemize}
  \item[（1）] 对任意的 \(\varepsilon>0\) ，在 \([a, b]\) 上使 \(\left|\lim _{t \rightarrow x} f(t)-f(x)\right|>\varepsilon\) 的点 \(x\) 至多有有限个．
  \item[（2）] \(f\) 在 \([a, b]\) 上至多只有可列个间断点．
\end{itemize}
\end{document}
